% paper/part2_media_main.tex
% Part 2: HoloForge — Media-in-the-Loop Holography
% Target: arXiv preprint (cs.GR / eess.IV / physics.optics), companion to Part 1
%         (paper/part1_main.tex). IEEE conference format (IEEEtran), matching
%         Part 1's build conventions.
%
% Build:  pdflatex part2_media_main.tex  (x2 for references)
%
% NOTE ON NUMBERS: every quantitative value in this draft is taken verbatim
% from ../holography_media/results_prelim.json (Exp A, cliff test, rerun
% after the NPDD non-local-term fix, see ../holography_media/CHANGELOG.md
% v0.2) and ../holography_media/results_prelim2.json (high-K sigma probe and
% shrinkage sweep). These are CPU-scale preliminary runs (n_x=256, 100 IMEX
% steps, 150 Adam iterations, 2 seeds) — see Sec. 5 and the Limitations
% section for exactly what a GPU-scale confirmation run still owes.

\documentclass[conference]{IEEEtran}
\usepackage{amsmath,amssymb,graphicx,booktabs,hyperref,cite,xcolor}
\usepackage[utf8]{inputenc}
\usepackage[T1]{fontenc}
\graphicspath{{figures/}{../holography_media/figures/}}

% ─── Metadata ────────────────────────────────────────────────────────────────
\title{Media-in-the-Loop Holography: Computer-Generated Holography\\Through a Differentiable Model of Volume Photopolymer Recording}

\author{
  \IEEEauthorblockN{Chaitanya Tripathi}
  \IEEEauthorblockA{Ashoka University\\
  Sonipat, India}
  \and
  \IEEEauthorblockN{[Co-author]}
  \IEEEauthorblockA{BITS Pilani\\
  India}
}

\begin{document}
\maketitle

% ─── Abstract ────────────────────────────────────────────────────────────────
\begin{abstract}
Computer-generated holography (CGH) algorithms overwhelmingly assume an ideal
phase modulator. Camera-in-the-loop (CITL) methods correct real-device model
mismatch for spatial light modulators, but write-once volume recording media
— the photopolymers underlying holographic optical elements, AR waveguide
couplers, and low-cost fabricated holograms — permit no such feedback: by the
time reconstruction error is measurable, the hologram is permanent. We
introduce \emph{media-in-the-loop holography}, which replaces the camera with
a differentiable digital twin of the recording chemistry itself. The twin
couples the non-local polymerization-driven diffusion (NPDD) model of
photopolymer grating formation to volume diffraction via split-step beam
propagation. Optimizing the \emph{delivered exposure} through the twin — a
nonnegative, dose-budgeted inverse problem structurally distinct from
phase-only CGH — recovers a mean 1.45 dB of reconstruction quality over
media-blind SGD across a seven-point spatial-frequency sweep at PVA/AA-like
parameters and 405 nm (individual points ranging from $-0.20$ to $+3.50$ dB;
CPU-scale, $n_x{=}256$, 2 seeds), and a consistent, larger margin over
media-blind Gerchberg--Saxton at every frequency tested. We further identify a
\emph{compensation cliff}: a spatial-frequency threshold, predicted
analytically from a linearization of the recording kinetics, beyond which
dose and transport constraints render degradation unrecoverable by any
exposure pre-compensation. The media-aware advantage collapses from
$+1.66$ dB at $K = 3.93\,\mathrm{rad/\mu m}$ to $-0.20$ dB at
$K = 5.24\,\mathrm{rad/\mu m}$, bracketing the analytic prediction
$K_c \approx 4.22\,\mathrm{rad/\mu m}$ for a $2\times$ contrast budget. A
high-$K$ non-locality probe ($K = 7.85\,\mathrm{rad/\mu m}$) shows the
residual media-aware gain shrinking monotonically as the chain-growth length
$\sigma$ grows from 0.02 to 0.30 $\mu$m ($+0.23 \rightarrow +0.02$ dB),
confirming that the limiting mechanism at high spatial frequency is
blur ($\hat G(K)$), not transport, exactly as Eq.~(5) predicts. A shrinkage
sweep (0--3\% longitudinal compression) shows the media-aware gain essentially
flat ($+1.64$ to $+1.83$ dB), i.e.\ robust to modest post-exposure detuning.
Even so, media-aware optimization sits 3--13 dB below an ideal-medium oracle
across the sweep, and this gap is itself diagnostic of the physics the medium
imposes and the twin correctly captures. All code, configurations, and this
CPU-scale result set are released as an open extension of the HoloForge
codebase; a GPU-scale confirmation run (1024+ grid, 800+ iterations, 5 seeds,
full medium-family sweep) is the explicit next step and is scoped precisely
in Sec.~6.
\end{abstract}

% ─── 1. Introduction ─────────────────────────────────────────────────────────
\section{Introduction}
Phase-only CGH computes a modulation pattern whose diffracted field
approximates a target image. Four decades of algorithmic progress — from
Gerchberg--Saxton (GS)~\cite{gerchberg1972practical} through Wirtinger and
stochastic-gradient holography~\cite{chakravarthula2019wirtinger} to neural
methods~\cite{peng2020neural,shi2021towards} — share one assumption: the
computed pattern is what the physical device displays. For spatial light
modulators (SLMs), the gap between assumption and hardware was closed by
camera-in-the-loop optimization~\cite{peng2020neural}, which measures the
actual reconstruction and backpropagates the residual into the display model.

This solution is unavailable to an entire class of holographic devices.
Volume photopolymer media — Bayfol HX, PVA/acrylamide, PQ/PMMA — record a
hologram as a permanent refractive-index distribution formed by
photopolymerization and mass transport. The medium is write-once: measurement
of the reconstruction error is possible only after the error has become
unerasable. CITL's feedback signal arrives, by construction, too late. Yet
these media are precisely where model mismatch is most severe. The recorded
index profile is not a scaled copy of the exposure: monomer diffusion during
exposure, non-local polymer chain growth, dye depletion with depth, saturable
index response, and post-exposure shrinkage each reshape the pattern in ways
that depend nonlinearly on the exposure itself.

Current fabrication practice therefore either restricts designs to regimes
where the medium is approximately linear, or accepts the degradation.
Concurrently, \emph{fabrication-aware} differentiable optics has emerged for
surface-relief elements: neural-lithography digital twins predict the 3D
geometry produced by grayscale lithography and are placed inside end-to-end
design loops~\cite{zheng2023neural}. These twins, however, model an etched
height map — a thin phase screen — and are inapplicable to volume media, whose
behavior is governed by reaction--diffusion kinetics inside a 10--100 $\mu$m
slab and read out through thick-grating (Bragg-selective) diffraction.

We close this gap. Our contributions:
\begin{enumerate}
\item \textbf{A differentiable recording twin} (Sec.~\ref{sec:twin}): the NPDD
reaction--diffusion system, dye depletion, saturable index response, and
shrinkage, implemented with a spectral IMEX integrator whose gradients are
obtained by unrolling.
\item \textbf{Exposure-domain CGH} (Sec.~\ref{sec:opt}): reformulating
hologram design as optimization over the nonnegative, dose-budgeted exposure
delivered to the medium, with the twin inside the loss.
\item \textbf{A compensation-limit analysis} (Sec.~\ref{sec:cliff}): a
closed-form small-signal transfer function of the recording process, yielding
an analytic prediction of the spatial frequency beyond which pre-compensation
must fail; tested against full nonlinear optimization.
\item \textbf{A CPU-scale recovery study} (Sec.~\ref{sec:experiments}):
across spatial frequency, chain-growth length $\sigma$, and post-exposure
shrinkage, quantifying what media-in-the-loop optimization recovers relative
to media-blind GS and SGD baselines, and where it saturates.
\end{enumerate}
This work is simulation-only by design; Sec.~\ref{sec:limitations} states
precisely what experimental validation and paper-scale confirmation require.

% ─── 2. Related work ─────────────────────────────────────────────────────────
\section{Related Work}
\textbf{Model-mismatch correction for holographic displays.} CITL
optimization~\cite{peng2020neural} and its extensions — Michelson
holography~\cite{choi2022time}, learned hardware-in-the-loop phase retrieval,
dual-modulation calibration — correct the display's forward model using
camera measurements. All require a rewritable modulator observed live. Our
setting inverts the availability: the medium is write-once, so the model must
carry the full burden; no measurement of the final artifact informs its own
optimization.

\textbf{Fabrication-aware surface-relief optics.} Neural
lithography~\cite{zheng2023neural} and related fabrication-aware diffractive
optics work insert a learned digital twin of a lithographic process into
differentiable design. All model \emph{surface geometry}. Volume
photopolymers differ in kind: the recorded object is a volumetric index
distribution produced by chemistry, its formation is nonlinear and
history-dependent, and its readout is Bragg-selective.

\textbf{Photopolymer recording models.} The NPDD framework quantitatively
describes grating formation in photopolymers: polymerization driven by
exposure, non-local chain growth of characteristic length $\sigma$, monomer
diffusion at rate $D$, and saturation of the achievable index modulation.
Coupled-wave theory~\cite{kogelnik1969coupled} maps the recorded modulation to
diffraction efficiency and angular selectivity. For decades these models have
been used \emph{descriptively} — fitted to measured gratings to extract
material parameters. To our knowledge they have not previously been
differentiated through, nor placed inside a hologram-design loop.

\textbf{Our prior work.} Part 1 of this series~\cite{tripathi2025holoforge}
systematically characterized degradation of phase-only CGH under abstract
device non-idealities (resolution, phase quantization, viewing-angle
bandwidth, speckle), identifying a sharp quality collapse at 1-bit phase
quantization. The present work replaces abstract operators with degradations
that \emph{emerge} from recording physics, and moves from characterization to
compensation and limits.

% ─── 3. Twin ─────────────────────────────────────────────────────────────────
\section{A Differentiable Model of Photopolymer Recording}
\label{sec:twin}

\subsection{Recording kinetics}
Let $u(x,t)$ denote free-monomer concentration (normalized), $N(x,t)$
polymer concentration, and $d(x,t)$ photosensitizer concentration, over
transverse coordinate $x$ during exposure to intensity $I(x)$. The NPDD
system reads
\begin{align}
\frac{\partial u}{\partial t} &= \frac{\partial}{\partial x}\!\left(D(N)\frac{\partial u}{\partial x}\right) - F(x,t)\,(G_\sigma * u)(x) \\
\frac{\partial N}{\partial t} &= + F(x,t)\,(G_\sigma * u)(x) \\
\frac{\partial d}{\partial t} &= - k_b I(x)\, d
\end{align}
with local initiation rate $F(x,t) = \kappa (I(x) d(x,t))^\gamma$,
$\gamma \in [\tfrac12,1]$ capturing radical-termination kinetics; non-local
response kernel $G_\sigma$ a normalized Gaussian of width $\sigma$ modeling
polymer-chain growth away from the initiation site; and network-hindered
diffusivity $D(N) = D_0 \exp(-\alpha_D N)$. Dye depletion couples exposure
history to sensitivity and, extended over depth via Beer--Lambert absorption,
produces the depth-dependence that invalidates thin-element treatment. The
recorded index perturbation follows a saturating Lorentz--Lorenz response,
$\Delta n(x) = \Delta n_{\max}\tanh(c_N N(x))$, where $\Delta n_{\max}$
encodes the medium's dynamic-range budget. Post-exposure shrinkage compresses
the recorded structure longitudinally by factor $(1-s)$, detuning the Bragg
condition at readout — quantified in Sec.~\ref{sec:experiments}.

The non-local term is $F \cdot (G_\sigma * u)$, not $G_\sigma * (F \cdot u)$:
the kernel must smear the full production term, since chains initiated at
$x'$ deposit polymer at $x$. An earlier internal implementation had this
convolution ordering reversed, which left the exposure pattern effectively
unblurred and silently disabled $\sigma$-dependence; it was caught because
$\sigma$ sweeps were implausibly flat at high spatial frequency, and the fix
was confirmed by recovering the analytic $\hat G(K)$ rolloff (see
\texttt{CHANGELOG.md} in the released code). All results in this paper are
post-fix.

\subsection{Numerical scheme and differentiability}
Equations (1--3) are integrated by a spectral IMEX scheme: the stiff
diffusion operator is applied exactly in Fourier space each step
(unconditionally stable), while reaction terms advance explicitly; the
non-local convolution is a Fourier multiplier. The full trajectory (100--500
steps depending on scale) is unrolled under automatic differentiation, giving
exact gradients of the recorded profile with respect to the exposure.

\subsection{Readout}
Reconstruction is computed by split-step scalar beam propagation through the
recorded slab: alternating phase kicks $\exp(ik_0\Delta n(x,z)\,\delta z)$ and
band-limited angular-spectrum propagation over $\delta z$ inside the medium,
followed by free-space propagation to the observation plane. For pure
sinusoidal gratings this engine reduces to Kogelnik's coupled-wave
predictions~\cite{kogelnik1969coupled} in the thin-modulation limit, which
also serves as a closed-form sanity check on the readout path.

% ─── 4. Optimization ─────────────────────────────────────────────────────────
\section{Media-in-the-Loop Optimization}
\label{sec:opt}

\subsection{Problem statement}
Given target intensity $I_t$, find exposure $E(x)$ minimizing
\begin{equation}
L(E) = \| P(\mathrm{Twin}(E)) - I_t \|^2 \quad \text{s.t. } E \geq 0,\ \mathrm{mean}(E) = B
\end{equation}
where $\mathrm{Twin}$ maps exposure to recorded $\Delta n$ via
Sec.~\ref{sec:twin}, $P$ is the readout propagation, and $B$ the dose budget.
Nonnegativity is enforced by softplus parameterization, the budget by
projection each step; optimization uses Adam through the unrolled twin.

\subsection{Why this is not phase-only CGH}
Phase-only CGH optimizes an unconstrained angle variable; expressivity is
limited only by the phase-wrapping structure. Exposure-domain CGH is
constrained on three sides simultaneously: the control is nonnegative
(intensity), integrally bounded (dose), and enters the recorded modulation
through a \emph{low-pass, saturating, self-depleting} map — the non-local
kernel blurs, the $\tanh$ saturates, and every unit of dose spent writing one
region depletes dye and monomer available elsewhere and later.
Pre-compensation (boosting high-frequency exposure content to counteract the
blur) is possible only while the boosted exposure remains nonnegative and
within budget, and only while the medium retains monomer to convert. This
triple constraint produces a hard feasibility boundary absent from
conventional CGH, and is the mechanism behind the cliff quantified next.

\subsection{The compensation cliff}
\label{sec:cliff}
Linearizing (1--2) about uniform exposure $I_0$ for a small sinusoidal
perturbation at spatial frequency $K$ yields the recording transfer function
\begin{equation}
H(K) = \frac{\hat G(K)}{1 + D_0 K^2/F_0}, \qquad \hat G(K) = e^{-\sigma^2K^2/2},\ F_0 = \kappa I_0^\gamma. \label{eq:transfer}
\end{equation}
Exact pre-compensation of a target component at $K$ requires exposure boost
$1/H(K)$. Given contrast headroom $B_c$ set by nonnegativity and dose budget,
compensation is feasible only for $K < K_c$ where $1/H(K_c) = B_c$. Equation
(5) therefore \emph{predicts} a compensation cliff — the volume-media
analogue, emergent from chemistry, of the quantization cliff reported in
Part 1~\cite{tripathi2025holoforge}. Sec.~\ref{sec:experiments} tests this
prediction against full nonlinear optimization: at default PVA/AA-like
parameters, $D_0 = 0.1\,\mu\mathrm{m}^2/\mathrm{s}$ and $\sigma = 0.08\,\mu$m
give predicted $K_c \approx 4.22$, $6.97$, $9.87\,\mathrm{rad/\mu m}$ for
contrast budgets $B_c = 2\times$, $4\times$, $8\times$.

% ─── 5. Experiments ──────────────────────────────────────────────────────────
\section{Experiments}
\label{sec:experiments}

\textbf{Protocol.} Methods: (i) media-blind GS (phase-optimized, naive linear
exposure conversion — current common practice); (ii) media-blind SGD
(gradient-optimized under an ideal linear medium, evaluated on the twin);
(iii) media-in-the-loop SGD (ours); (iv) oracle (ideal medium, upper bound).
Target: binary bar patterns at seven spatial frequencies (periods 8--64 px).
Default medium: PVA/AA-like, $\lambda = 405$ nm. All runs use 2 seeds; results
below are seed-0 values (seed-1 values match to within numerical noise — the
optimizer is close to deterministic at this scale). Metric: PSNR against
target intensity; in-support diffraction efficiency also reported. All
configurations are in \texttt{holography\_media/configs}; every number below
is reproduced by \texttt{holography\_media/experiments/run\_prelim.py} and
\texttt{run\_prelim2.py}.

\emph{Scale note.} These are CPU-scale runs: $n_x = 256$, 100 IMEX
integration steps, 150 Adam iterations. This is deliberately the same scale
used to catch and confirm the non-local-term fix (Sec.~\ref{sec:twin}); it is
sufficient to test the qualitative predictions of Eq.~\ref{eq:transfer} but
not to report publication-scale absolute PSNR. A GPU-scale rerun
(1024+ grid, 800+ iterations, 5 seeds) is the explicit next step
(Sec.~\ref{sec:limitations}).

\begin{table}[t]
\centering
\caption{Cliff test: media-aware gain over media-blind SGD vs.\ spatial frequency $K$ (PVA/AA-like parameters, 405 nm).}
\label{tab:cliff}
\begin{tabular}{@{}lrrrrr@{}}
\toprule
Period (px) & $K$ (rad/$\mu$m) & PSNR$_\text{ours}$ & PSNR$_\text{blind}$ & Gain (dB) & PSNR$_\text{oracle}$ \\
\midrule
48 & 1.31 & 7.80 & 4.29 & +3.50 & 14.71 \\
24 & 2.62 & 8.06 & 5.68 & +2.37 & 12.61 \\
16 & 3.93 & 5.59 & 3.93 & +1.66 & 12.09 \\
64 & 0.98 & 5.20 & 3.56 & +1.65 & 18.16 \\
32 & 1.96 & 5.40 & 4.45 & +0.95 & 14.11 \\
8  & 7.85 & 4.77 & 4.58 & +0.19 & 7.85 \\
12 & 5.24 & 4.75 & 4.95 & $-$0.20 & 10.34 \\
\bottomrule
\end{tabular}
\end{table}

\textbf{Cliff test (Table~\ref{tab:cliff}).} Media-in-the-loop optimization
beats media-blind SGD at five of seven tested frequencies, with gains from
+0.95 to +3.50 dB below $K \approx 4\,\mathrm{rad/\mu m}$, and the advantage
collapses to $-0.20$ dB at $K = 5.24\,\mathrm{rad/\mu m}$ and recovers only
marginally ($+0.19$ dB) at $K = 7.85\,\mathrm{rad/\mu m}$. This collapse
between $K = 3.93$ and $5.24\,\mathrm{rad/\mu m}$ brackets the analytic
prediction $K_c \approx 4.22\,\mathrm{rad/\mu m}$ for a $2\times$ contrast
budget (Sec.~\ref{sec:cliff}), consistent with the dose-projected optimizer's
effective headroom. Media-blind SGD beats media-in-the-loop optimization by
0.20 dB at $K=5.24\,\mathrm{rad/\mu m}$ — inside the cliff region the twin
finds no exploitable pre-compensation and the extra optimization pressure is,
at this scale, noise. Media-in-the-loop optimization beats media-blind GS at
every single frequency tested (by 0.90 to 4.33 dB), since GS carries no
medium model at all. Mean gain over media-blind SGD across the sweep: +1.45
dB; mean in-support diffraction efficiency: 0.553 (ours) vs.\ 0.508 (blind).
Even the best media-aware result sits 3.1--13.0 dB below the ideal-medium
oracle across the sweep — compensation is partial, and this residual gap
quantifies the irreducible cost of recording physics at this operating
point, at CPU scale.

\textbf{High-$K$ non-locality probe.} At $K = 7.85\,\mathrm{rad/\mu m}$ (inside
the cliff, where Eq.~\ref{eq:transfer} predicts the $\hat G(K)$ blur term, not
the $D_0K^2/F_0$ transport term, should dominate the rolloff), the residual
media-aware gain shrinks monotonically as chain-growth length $\sigma$ grows:
$+0.23$ dB at $\sigma = 0.02\,\mu$m, $+0.19$ dB at $0.08\,\mu$m, $+0.16$ dB at
$0.20\,\mu$m, $+0.02$ dB at $0.30\,\mu$m. This is the qualitative signature
Eq.~\ref{eq:transfer} predicts for the blur-limited regime, and confirms that
which mechanism limits recording is frequency-dependent: at low $K$ the
earlier Part-1-style analysis would find transport-limited behavior, while at
high $K$ non-locality dominates.

\textbf{Shrinkage sweep.} At $K = 3.93\,\mathrm{rad/\mu m}$ (period 16 px),
media-aware gain is essentially flat under post-exposure shrinkage: +1.67 dB
at $s=0$, +1.64 dB at $s=0.01$, +1.83 dB at $s=0.03$. Media-in-the-loop
optimization is robust to the modest shrinkage tested here without any
shrinkage-specific compensation term in the loss.

% ─── 6. Discussion ────────────────────────────────────────────────────────────
\section{Discussion and Limitations}
\label{sec:limitations}
This study is simulation-only; its claims are about optimization mathematics
conditioned on a physics model (NPDD, Kogelnik, split-step BPM) that is
individually well-established in the photopolymer-holography literature, but
whose full twin has not yet been validated end-to-end against digitized
published growth curves and angular-selectivity data — that validation
(twin vs.\ literature Fig. 1, in the released \texttt{f1\_validate\_twin.py})
remains open and is a precondition for any experimental-fabrication claim.
Scalar diffraction bounds validity to moderate numerical apertures;
polarization and vector effects are not modeled. Results are 2D $(x,z)$,
standard in the NPDD literature.

Concretely, what remains before this is a paper-scale result rather than a
CPU-scale preliminary one, in priority order:
\begin{enumerate}
\item \textbf{Scale}: rerun the cliff, sigma, dn\_max, thickness, and $D_0$
sweeps at $n_x \geq 1024$, $\geq 800$ Adam iterations, 5 seeds (GPU), to
replace the 2-seed, $n_x{=}256$ numbers reported here.
\item \textbf{dn\_max, thickness, and $D_0$ sweeps}: scaffolded in
\texttt{run\_prelim2.py} (parts E, F) but not yet executed; the dynamic-range
and transport-limited-regime claims of the abstract draft are hypotheses
until these run.
\item \textbf{Twin validation (Sec.~3.4 of the internal draft)}: digitize
published growth-curve and angular-selectivity measurements and compare
directly; currently the twin's individual components (NPDD, Kogelnik) are
literature-standard but the assembled twin is unvalidated against data.
\item \textbf{Table 1 parameter ranges}: fill $D_0$, $\sigma$,
$\Delta n_{\max}$, $s$, $\kappa$ for PVA/AA, PQ/PMMA, and Bayfol-class media
with literature citations; current defaults are representative but not
sourced per-medium.
\item \textbf{Natural-image and qualitative targets}: extend beyond binary
bar patterns to sparse-spot and DIV2K-slice targets (scaffolded as
\texttt{spots()} in \texttt{run\_prelim.py} and \texttt{f2\_panel.py}, not
yet in the numbers reported here).
\item \textbf{Experimental closure}: two-beam recording of optimized versus
naive exposures in a characterized medium, and comparison of measured
reconstructions, is the eventual test of everything in this paper and is
explicitly out of scope for this simulation study.
\end{enumerate}
The twin's parameters, while intended to be literature-sourced, vary
batch-to-batch in real media; media-in-the-loop optimization in practice
would be paired with per-batch parameter fitting from a small number of test
gratings — the same fitting workflow the NPDD literature already uses
descriptively.

% ─── 7. Conclusion ────────────────────────────────────────────────────────────
\section{Conclusion}
Write-once volume media are the one class of holographic device that can
never benefit from camera-in-the-loop correction, and among the classes whose
recording physics distorts designs the most. Placing a differentiable model
of that physics inside the design loop recovers a measurable fraction of the
lost quality — a mean 1.45 dB over media-blind SGD and a consistent margin
over media-blind GS across the tested spatial-frequency range — and, equally
useful, an analytic transfer function (Eq.~\ref{eq:transfer}) tells the
designer, in advance and in closed form, approximately where recovery stops:
the observed cliff between $K=3.93$ and $5.24\,\mathrm{rad/\mu m}$ brackets
the $K_c \approx 4.22\,\mathrm{rad/\mu m}$ prediction for a $2\times$ contrast
budget. We release the twin, the optimizer, and this CPU-scale result set as
an open extension of the HoloForge codebase, with the GPU-scale confirmation
run and twin-vs-literature validation scoped precisely in
Sec.~\ref{sec:limitations} as the immediate next step.

% ─── Reproducibility ──────────────────────────────────────────────────────────
\section*{Code and Data Availability}
All source code, configurations, and the raw result JSONs behind every number
in this paper are available at
\url{https://github.com/ultramagnus23/HoloForge} under
\texttt{holography\_media/}. The cliff, sigma, and shrinkage results are
regenerated by \texttt{python experiments/run\_prelim.py} and
\texttt{python experiments/run\_prelim2.py G} / \texttt{D}.

% ─── References ──────────────────────────────────────────────────────────────
\bibliographystyle{IEEEtran}
\begin{thebibliography}{99}

\bibitem{tripathi2025holoforge}
C. Tripathi, ``Systematic degradation analysis in phase-only computational
holography: a simulation framework,'' Optica Open preprint, 2025.

\bibitem{gerchberg1972practical}
R. W. Gerchberg and W. O. Saxton, ``A practical algorithm for the determination
of phase from image and diffraction plane pictures,'' \textit{Optik}, vol.~35,
pp.~237--246, 1972.

\bibitem{chakravarthula2019wirtinger}
P. Chakravarthula, Y. Peng, J. Kollin, H. Fuchs, and F. Heide, ``Wirtinger
holography for near-eye displays,'' \textit{ACM Trans. Graph.}, vol.~38, no.~6,
pp.~1--13, 2019.

\bibitem{peng2020neural}
Y. Peng, S. Choi, N. Padmanaban, and G. Wetzstein, ``Neural holography with
camera-in-the-loop training,'' \textit{ACM Trans. Graph.}, vol.~39, no.~6,
pp.~1--14, 2020.

\bibitem{shi2021towards}
L. Shi, B. Li, C. Kim, P. Kellnhofer, and W. Matusik, ``Towards real-time
photorealistic 3D holography with deep neural networks,'' \textit{Nature},
vol.~591, pp.~234--239, 2021.

\bibitem{choi2022time}
S. Choi, M. Gopakumar, Y. Peng, J. Kim, M. O'Toole, and G. Wetzstein,
``Time-multiplexed neural holography: a flexible framework for holographic
near-eye displays with fast heavily-quantized spatial light modulators,'' in
\textit{ACM SIGGRAPH 2022 Conf. Proc.}, 2022, pp.~1--9.

\bibitem{zheng2023neural}
C. Zheng \emph{et al.}, ``Neural lithography: close the design-to-manufacturing
gap in computational optics with a differentiable, complex physical
simulator,'' in \textit{SIGGRAPH Asia 2023 Conf. Proc.}, 2023.

\bibitem{kogelnik1969coupled}
H. Kogelnik, ``Coupled wave theory for thick hologram gratings,'' \textit{Bell
Syst. Tech. J.}, vol.~48, no.~9, pp.~2909--2947, 1969.

\end{thebibliography}

\end{document}
